\documentclass[11pt]{article}
\usepackage[margin=1in]{geometry}
\usepackage{amsmath,amssymb,amsthm}
\usepackage{booktabs}
\usepackage{hyperref}

\newtheorem{theorem}{Theorem}
\newtheorem{lemma}[theorem]{Lemma}
\newtheorem{proposition}[theorem]{Proposition}
\newtheorem{corollary}[theorem]{Corollary}
\theoremstyle{definition}
\newtheorem{definition}[theorem]{Definition}
\theoremstyle{remark}
\newtheorem{remark}[theorem]{Remark}

\newcommand{\wt}{\operatorname{wt}}
\newcommand{\ch}{\operatorname{ch}}
\newcommand{\sgn}{\operatorname{sgn}}
\newcommand{\Z}{\mathbb Z}
\newcommand{\tb}[2]{\mathtt{#1/#2}}

\title{Disproof of Conjecture 00000003831:\\
a seminormal $\mathfrak{gl}_3$ crystal whose character is $s_{21}-s_{111}$}
\author{jilint777}
\date{2026-10-04}

\begin{document}
\sloppy
\maketitle

\begin{abstract}
Conjecture 00000003831 asserts that for every finite crystal $B$ the Schur expansion of the
Weyl symmetrization of the character $\ch(B)$ has nonnegative coefficients, and that for
$B=B(\lambda)\otimes B(\mu)$ these coefficients are the Littlewood--Richardson coefficients.
The first clause is false. We exhibit a connected $7$-vertex \emph{seminormal} $\mathfrak{gl}_3$
crystal $B_7$, obtained from $B(2,1,0)$ by identifying its two vertices of weight $(1,1,1)$
into one vertex and keeping all eight edges. It satisfies all of Kashiwara's crystal axioms,
$\varepsilon_i,\varphi_i$ are the string lengths, and
\[
\ch(B_7)=m_{21}+m_{111}=s_{21}-s_{111}.
\]
The character is symmetric, so the Weyl symmetrizer $J(x^\rho f)/J(x^\rho)$ fixes it and the
coefficient of $s_{111}$ is $-1$; the orbit sum gives $6s_{21}-6s_{111}$ and the orbit average
gives $s_{21}-s_{111}$. Seven vertices is the minimum for a seminormal $\mathfrak{gl}_3$
counterexample, and seminormal $\mathfrak{gl}_2$ crystals never give one. Clause~1 is true for
\emph{normal} (Stembridge) crystals, so it can only be rescued by normal-type readings of
``finite crystal'' (normal crystals, or the affine notion of a finite crystal, which is
classically normal). The refutation is formalised in Lean~4 (core library only): crystals,
Kashiwara's axioms, seminormality, characters, Schur polynomials (from tableaux), the
antisymmetrizer and the Weyl symmetrization are defined from scratch. An independent Python
script rechecks everything and performs the minimality searches.
\end{abstract}

\section{The claim}

The conjecture (English version): \emph{``Definition: the crystal character $\ch(B)$ denotes the
generating function of vertex weights. Conjecture: for any finite crystal, the coefficients of the
Schur expansion of the Weyl symmetrization of $\ch(B)$ are nonnegative, and for
$B=B(\lambda)\otimes B(\mu)$ these coefficients are exactly the classical Littlewood--Richardson
coefficients.''}

It is a conjunction of two clauses:
\begin{itemize}
\item[(C1)] for every finite crystal $B$, the Schur expansion of the Weyl symmetrization of
$\ch(B)$ has nonnegative coefficients;
\item[(C2)] for $B=B(\lambda)\otimes B(\mu)$ these coefficients are the LR coefficients
$c^\nu_{\lambda\mu}$.
\end{itemize}
We refute (C1); one false clause makes the conjunction false. Under the Weyl-symmetrizer reading
(C2) is \emph{true}: $\ch(B(\lambda)\otimes B(\mu))=s_\lambda s_\mu$ is symmetric, so it is fixed
by the symmetrizer, and $s_\lambda s_\mu=\sum_\nu c^\nu_{\lambda\mu}s_\nu$. The disproof rests
entirely on (C1).

\section{Definitions}

\paragraph{Crystals.} We work with $\mathfrak{gl}_3$: weight lattice $P=\Z^3$ with basis
$\epsilon_1,\epsilon_2,\epsilon_3$, simple roots $\alpha_1=\epsilon_1-\epsilon_2$,
$\alpha_2=\epsilon_2-\epsilon_3$, and coroots acting by $\langle h_1,w\rangle=w_1-w_2$,
$\langle h_2,w\rangle=w_2-w_3$. Following Kashiwara
\cite{Kas93,Kas95} and Bump--Schilling \cite[Ch.~2]{BS17}, a \emph{crystal} is a set $B$ with
maps $\wt\colon B\to P$, $\varepsilon_i,\varphi_i\colon B\to\Z\sqcup\{-\infty\}$ and
$\tilde e_i,\tilde f_i\colon B\to B\sqcup\{0\}$ ($i=1,2$) such that
\begin{enumerate}
\item $\varphi_i(b)=\varepsilon_i(b)+\langle h_i,\wt(b)\rangle$;
\item $\tilde f_i b=b'$ if and only if $\tilde e_i b'=b$;
\item if $\tilde e_ib\in B$ then $\wt(\tilde e_ib)=\wt(b)+\alpha_i$,
$\varepsilon_i(\tilde e_ib)=\varepsilon_i(b)-1$, $\varphi_i(\tilde e_ib)=\varphi_i(b)+1$;
\item if $\tilde f_ib\in B$ then $\wt(\tilde f_ib)=\wt(b)-\alpha_i$,
$\varepsilon_i(\tilde f_ib)=\varepsilon_i(b)+1$, $\varphi_i(\tilde f_ib)=\varphi_i(b)-1$;
\item if $\varphi_i(b)=-\infty$ then $\tilde e_ib=\tilde f_ib=0$.
\end{enumerate}
A crystal is \emph{seminormal} if
$\varepsilon_i(b)=\max\{k\ge 0:\tilde e_i^k b\ne 0\}$ and
$\varphi_i(b)=\max\{k\ge0:\tilde f_i^kb\ne0\}$ for all $i,b$; it is \emph{normal} if it is
(isomorphic to) the crystal of an integrable $U_q(\mathfrak{gl}_3)$-module, equivalently a
disjoint union of highest-weight crystals $B(\lambda)$ (for simply-laced types this is
equivalent to Stembridge's local axioms \cite{Ste03}). In these references seminormality and
normality are \emph{extra} conditions; ``crystal'' alone means the abstract notion above.

\paragraph{Characters and Schur polynomials.} For a finite crystal,
$\ch(B)=\sum_{b\in B}x^{\wt(b)}\in\Z[x_1^{\pm1},x_2^{\pm1},x_3^{\pm1}]$. The Weyl group $S_3$
permutes the variables. Let $J(f)=\sum_{w\in S_3}\sgn(w)\,w(f)$, $\rho=(2,1,0)$ and
$a_\mu=J(x^\mu)$. For a dominant weight $\lambda$ ($\lambda_1\ge\lambda_2\ge\lambda_3$, integers)
the Schur polynomial is $s_\lambda=a_{\lambda+\rho}/a_\rho$; for a partition $\lambda$ it equals
$\sum_T x^T$ over semistandard Young tableaux of shape $\lambda$ with entries in $\{1,2,3\}$,
and $s_{\lambda+(k,k,k)}=(x_1x_2x_3)^ks_\lambda$.

\begin{lemma}[uniqueness of Schur expansions]\label{lem:unique}
The $s_\lambda$, $\lambda$ dominant, are linearly independent, and a symmetric Laurent
polynomial $f$ equals $\sum_\lambda c_\lambda s_\lambda$ with
$c_\lambda=[x^{\lambda+\rho}]\,J(x^\rho f)$.
\end{lemma}
\begin{proof}
If $f$ is symmetric then $J(x^\rho f)=f\,a_\rho$, an antisymmetric Laurent polynomial. An
antisymmetric Laurent polynomial is determined by its coefficients on the strictly decreasing
exponents $\mu_1>\mu_2>\mu_3$ (every exponent with two equal entries has coefficient $0$, and
the others are $\pm$ a strictly decreasing one), so $J(x^\rho f)=\sum_\lambda c_\lambda
a_{\lambda+\rho}$ with $c_\lambda$ as stated. Dividing by $a_\rho$ gives the expansion. If
$\sum c_\lambda s_\lambda=0$ then $\sum c_\lambda a_{\lambda+\rho}=0$, and the $a_\mu$ with
distinct strictly decreasing $\mu$ have disjoint supports, so all $c_\lambda=0$.
\end{proof}

\paragraph{Readings of ``Weyl symmetrization''.} For a Laurent polynomial $f$:
\begin{itemize}
\item[(W)] the Weyl symmetrizer $\pi(f)=J(x^\rho f)/J(x^\rho)$ (the Weyl character formula
operator, equal to the Demazure operator $\pi_{w_0}=\pi_1\pi_2\pi_1$); its Schur coefficients are
$[x^{\lambda+\rho}]J(x^\rho f)$ for any $f$;
\item[(O)] the orbit sum $\sum_{w\in S_3}w(f)$;
\item[(A)] the orbit average $\frac16\sum_{w\in S_3}w(f)$.
\end{itemize}
If $f$ is symmetric then $\pi(f)=f$ (because $J(x^\rho f)=f\,J(x^\rho)$), the orbit sum is $6f$
and the average is $f$.

\section{The counterexample}

The crystal $B_7$ has seven vertices
\[
A\,(2,1,0),\ B\,(1,2,0),\ C\,(2,0,1),\ D\,(0,2,1),\ E\,(1,0,2),\ F\,(0,1,2),\ Z\,(1,1,1)
\]
(weights in brackets), and the edges
\[
\tilde f_1:\ A\to B,\quad C\to Z\to D,\quad E\to F;\qquad
\tilde f_2:\ A\to C,\quad B\to Z\to E,\quad D\to F,
\]
with $\tilde e_i$ the inverse partial maps and $\tilde e_i,\tilde f_i=0$ elsewhere. The values of
$\varepsilon_i,\varphi_i$ are the string lengths:
\begin{center}
\begin{tabular}{lccccccc}
\toprule
 & $A$ & $B$ & $C$ & $D$ & $E$ & $F$ & $Z$\\
\midrule
$\wt$ & $(2,1,0)$ & $(1,2,0)$ & $(2,0,1)$ & $(0,2,1)$ & $(1,0,2)$ & $(0,1,2)$ & $(1,1,1)$\\
$(\varepsilon_1,\varphi_1)$ & $(0,1)$ & $(1,0)$ & $(0,2)$ & $(2,0)$ & $(0,1)$ & $(1,0)$ & $(1,1)$\\
$(\varepsilon_2,\varphi_2)$ & $(0,1)$ & $(0,2)$ & $(1,0)$ & $(0,1)$ & $(2,0)$ & $(1,0)$ & $(1,1)$\\
\bottomrule
\end{tabular}
\end{center}
The crystal $B(2,1,0)$ has the eight tableaux
$\tb{11}{2},\tb{12}{2},\tb{11}{3},\tb{12}{3},\tb{13}{2},\tb{13}{3},
\tb{22}{3},\tb{23}{3}$ (written first row/second row), with
$\tilde f_1$-edges $\tb{11}{2}\to\tb{12}{2}$, $\tb{11}{3}\to\tb{12}{3}\to\tb{22}{3}$,
$\tb{13}{3}\to\tb{23}{3}$ and $\tilde f_2$-edges $\tb{11}{2}\to\tb{11}{3}$,
$\tb{12}{2}\to\tb{13}{2}\to\tb{13}{3}$, $\tb{22}{3}\to\tb{23}{3}$. The two
vertices of weight $(1,1,1)$ are $\tb{12}{3}$ (on a $1$-string of length $2$, alone in
colour $2$) and $\tb{13}{2}$ (alone in colour $1$, on a $2$-string of length $2$).
Identifying them to a single vertex $Z$ keeps all eight edges and gives exactly $B_7$.

\begin{proposition}\label{prop:crystal}
$B_7$ is a connected seminormal crystal (with all $\varepsilon_i,\varphi_i$ finite).
\end{proposition}
\begin{proof}
Each $\tilde f_i$ is injective on its domain, so $\tilde e_i$ is a well-defined inverse and
axiom~2 holds. Each edge $u\to v$ of colour $i$ has $\wt(u)-\wt(v)=\alpha_i$ (check the four
edges of each colour, e.g.\ $C\to Z$: $(2,0,1)-(1,1,1)=(1,-1,0)=\alpha_1$, $Z\to E$:
$(1,1,1)-(1,0,2)=(0,1,-1)=\alpha_2$). The $i$-strings are $A\to B$, $C\to Z\to D$, $E\to F$
for $i=1$ and $A\to C$, $B\to Z\to E$, $D\to F$ for $i=2$; every vertex lies on exactly one
string of each colour. Setting $\varepsilon_i(b)$ (resp.\ $\varphi_i(b)$) to the number of steps
from $b$ to the top (resp.\ bottom) of its string gives the table above, makes $B_7$ seminormal
by definition, and gives axioms~3 and~4 automatically. Axiom~1 is a direct check: for a string
$b_0\to\dots\to b_\ell$ of colour $i$ we have
$\langle h_i,\wt(b_j)\rangle=\langle h_i,\wt(b_0)\rangle-2j$, and from the table
$\langle h_i,\wt(b_0)\rangle=\ell$ in all six strings (for instance
$\langle h_1,\wt(C)\rangle=2-0=2$, $\langle h_2,\wt(B)\rangle=2-0=2$,
$\langle h_1,\wt(A)\rangle=1$). Hence $\varphi_i(b_j)-\varepsilon_i(b_j)=(\ell-j)-j
=\langle h_i,\wt(b_j)\rangle$. Axiom~5 is void. Connectedness: every vertex is joined to $A$.
\end{proof}

\begin{proposition}\label{prop:char}
$\ch(B_7)=s_{(2,1,0)}-s_{(1,1,1)}$, and this is the unique Schur expansion.
\end{proposition}
\begin{proof}
The weights of $B_7$ are the six permutations of $(2,1,0)$ and $(1,1,1)$, each once, so
$\ch(B_7)=m_{21}+m_{111}$, which is symmetric. The semistandard tableaux of shape $(2,1)$ with
entries $\le 3$ are the eight listed above, so $s_{21}=m_{21}+2m_{111}$, and $s_{111}=x_1x_2x_3
=m_{111}$. Hence $\ch(B_7)=s_{21}-s_{111}$. Uniqueness is Lemma~\ref{lem:unique}; in degree $3$
it is also visible directly, since the coefficients of $x^{300},x^{210},x^{111}$ in
$s_{300},s_{210},s_{111}$ form the unitriangular matrix
$\left(\begin{smallmatrix}1&0&0\\1&1&0\\1&2&1\end{smallmatrix}\right)$.
\end{proof}

\begin{theorem}\label{thm:main}
Clause (C1) fails for $B_7$ under each of the readings (W), (O), (A):
\[
\pi(\ch B_7)=s_{21}-s_{111},\qquad \sum_{w}w(\ch B_7)=6s_{21}-6s_{111},\qquad
\tfrac16\sum_w w(\ch B_7)=s_{21}-s_{111}.
\]
The coefficient of $s_{111}$ is $-1$, $-6$, $-1$ respectively. Consequently the conjecture is
false.
\end{theorem}
\begin{proof}
$\ch B_7$ is symmetric, so by the remark after the readings the three symmetrizations are
$f,6f,f$ with $f=\ch B_7$; Proposition~\ref{prop:char} and Lemma~\ref{lem:unique} give the
expansions. Explicitly, $J(x^\rho\ch B_7)=a_{(4,2,0)}-a_{(3,2,1)}$, so
$[x^{(1,1,1)+\rho}]J(x^\rho\ch B_7)=-1$.
\end{proof}

\section{What is true, and the boundary of the counterexample}

\paragraph{Normal crystals.} If $B$ is normal then $B\cong\bigsqcup_k B(\lambda^{(k)})$ and
$\ch B=\sum_k s_{\lambda^{(k)}}$ (the character of $B(\lambda)$ is the irreducible character
$s_\lambda$), which is symmetric and Schur-nonnegative; all three readings then give
nonnegative coefficients. In particular (C1) holds for $B(\lambda)\otimes B(\mu)$, and (C2) holds
under (W) and (A); under (O) the coefficients are $6c^\nu_{\lambda\mu}$, so (C2) fails under (O).
Our $B_7$ is not normal: a normal crystal has a Schur-positive character. So the only readings
under which (C1) survives are normal-type ones: ``finite crystal'' $=$ ``normal crystal'', or the
affine-theory term discussed in Section~\ref{sec:readings}. The literal statement
says ``any finite crystal'', and in the standard references \cite{Kas93,Kas95,BS17}
``crystal'' does not include normality (or even seminormality); $B_7$ is moreover seminormal,
the strongest condition short of normality that is usually attached to the word.

\paragraph{Characters of seminormal crystals.}
\begin{lemma}\label{lem:semi}
Let $B$ be a finite seminormal $\mathfrak{gl}_r$ crystal and $m(\nu)$ the number of vertices of
weight $\nu$. Then $m$ is $S_r$-invariant, and for every $i$ and every $\nu$ with
$\langle h_i,\nu\rangle\ge0$ we have $m(\nu+\alpha_i)\le m(\nu)$ (the \emph{string condition}).
Conversely every finitely supported $m\ge0$ with these two properties is the weight
multiplicity function of a finite seminormal crystal.
\end{lemma}
\begin{proof}
Seminormality makes every vertex lie on a unique $i$-string $b_0\to\dots\to b_\ell$ with
$\varepsilon_i(b_0)=0$, $\varphi_i(b_0)=\ell$, so $\langle h_i,\wt b_0\rangle=\ell$ and the
weights $\wt b_0-j\alpha_i$, $0\le j\le\ell$, are permuted by the reflection $s_i$. Hence $m$ is
$s_i$-invariant for all $i$, i.e.\ $S_r$-invariant. If $\langle h_i,\nu\rangle=k\ge0$, a string
through $\nu+\alpha_i$ contains $s_i(\nu+\alpha_i)=\nu-(k+1)\alpha_i$ and is contiguous, so it
contains $\nu$; this gives $m(\nu+\alpha_i)\le m(\nu)$. Conversely, for each colour separately,
cut every line $\nu+\Z\alpha_i$ into strings by the usual layering (the multiplicities along the
line are symmetric about the reflecting point and weakly decrease away from it), and define
$\tilde f_i,\tilde e_i,\varepsilon_i,\varphi_i$ along these strings. The axioms only involve one
colour at a time, so the colours do not interact.
\end{proof}

\begin{proposition}[$\mathfrak{gl}_2$]
Every finite seminormal $\mathfrak{gl}_2$ (or $\mathfrak{sl}_2$) crystal has a Schur-positive
character; so no counterexample exists for $\mathfrak{gl}_2$ among seminormal crystals.
\end{proposition}
\begin{proof}
There is one colour; $B$ is the disjoint union of its strings, and a string with top weight
$(a,c)$ has $\ell=a-c\ge0$ and character $\sum_{j=0}^{\ell}x_1^{a-j}x_2^{c+j}=s_{(a,c)}$.
\end{proof}
For $\mathfrak{gl}_3$ each colour still decomposes $\ch B$ into $\mathfrak{sl}_2$-characters,
but the two decompositions need not come from a common $\mathfrak{gl}_3$-structure; $B_7$
exploits exactly this.

\begin{proposition}[minimality]\label{prop:min}
Every finite seminormal $\mathfrak{gl}_3$ crystal with at most $6$ vertices has a Schur-positive
character. Hence $7$ is the minimum number of vertices of a seminormal $\mathfrak{gl}_3$
counterexample to (C1).
\end{proposition}
\begin{proof}
Edges preserve the degree $w_1+w_2+w_3$, and the Schur expansion is compatible with the
decomposition by degree, so a minimal counterexample lives in one degree $d$; multiplying by
$(x_1x_2x_3)^k$ maps $s_\lambda$ to $s_{\lambda+(k,k,k)}$ and preserves everything, so we may
take $d\in\{0,1,2\}$. If a dominant $\lambda$ occurs, $p=\lambda_1-\lambda_2$, then by
Lemma~\ref{lem:semi} the $p+1$ distinct weights $\lambda-j\alpha_1$ ($0\le j\le p$) all occur,
so $p\le N-1$, and likewise $\lambda_2-\lambda_3\le N-1$, where $N$ is the number of vertices.
This leaves a finite list of multiplicity functions; \texttt{verify.py} enumerates all of them
for $N\le 6$, keeps those satisfying the conditions of Lemma~\ref{lem:semi} (necessary
conditions suffice here) and finds exactly $12$, namely
$k\,s_{000}$ ($1\le k\le6$), $s_{100}$, $2s_{100}$, $s_{110}$, $2s_{110}$, $s_{200}$ and
$s_{(1,1,-1)}$, all Schur-positive. For $N=7$ there is exactly one non-positive admissible
character up to the shift, $m_{(1,0,-1)}+m_{000}=s_{(1,0,-1)}-s_{000}$, which is
$\ch(B_7)$ shifted by $(x_1x_2x_3)^{-1}$; the script also builds a seminormal crystal for it
and rechecks the axioms.
\end{proof}

\begin{remark}[a cheaper abstract example]
Without seminormality, one vertex suffices: the $\mathfrak{gl}_2$ crystal $\{b\}$ with
$\wt(b)=(-1,2)$, $\varepsilon(b)=3$, $\varphi(b)=0$, $\tilde e b=\tilde f b=0$ satisfies
Kashiwara's axioms ($\varphi-\varepsilon=-3=\langle h,\wt b\rangle$). Then
$J(x^\rho x_1^{-1}x_2^2)=a_{(0,2)}=-a_{(2,0)}$, so the Weyl symmetrization is $-s_{(1,0)}$, and
the orbit sum is $s_{(2,-1)}-s_{(1,0)}$. A second one: the lowest vertex $b$ of $B(2,0)$, kept
alone with its inherited values $\wt(b)=(0,2)$, $\varepsilon(b)=2$, $\varphi(b)=0$ and no edges,
is a crystal ($\varphi-\varepsilon=-2=0-2$). Its character is $x_2^2$, and
$J(x^\rho x_2^2)=x_1x_2^2-x_1^2x_2=-a_{(2,1)}$, so the Weyl symmetrization is $-s_{(1,1)}$
(orbit sum $x_1^2+x_2^2=s_{(2,0)}-s_{(1,1)}$). We do not rely on these examples: it is far from the
spirit of the conjecture, while $B_7$ is seminormal, connected and differs from $B(2,1,0)$
only by one vertex.
\end{remark}

\section{Readings covered}\label{sec:readings}
\begin{itemize}
\item \emph{Crystal}: abstract Kashiwara crystals and seminormal crystals ($B_7$ is
seminormal). Only normal-type readings rescue (C1).
\item \emph{The symmetrization would be idle for normal crystals.} For normal (indeed
seminormal, Lemma~\ref{lem:semi}) crystals $\ch(B)$ is already $W$-invariant, so ``Weyl
symmetrization'' would be vacuous, and for normal crystals (C1) would be a triviality. The explicit
symmetrization therefore points to the abstract notion, under which even a $1$-vertex crystal is
a counterexample (the remark above). $B_7$ refutes (C1) even under the stronger seminormal
reading.
\item \emph{Affine reading.} In affine crystal theory, a ``finite crystal'' means the crystal
base of a finite-dimensional $U'_q(\hat{\mathfrak g})$-module (e.g.\ Kirillov--Reshetikhin or
perfect crystals). Such crystals are classically normal: their restriction to the classical
subalgebra is a disjoint union of $B(\lambda)$'s, so (C1) holds under this reading. It is one of
the normal-type readings named above, and our counterexample does not apply to it.
\item \emph{Weyl symmetrization}: the symmetrizer $J(x^\rho f)/J(x^\rho)$ (Demazure operator
$\pi_{w_0}$), its bialternant coefficients $[x^{\lambda+\rho}]J(x^\rho f)$, the orbit sum and
the orbit average all give a negative coefficient. Since $\ch B_7$ is symmetric, any
symmetrization that fixes symmetric functions, or multiplies them by a positive constant, does.
\item \emph{$\mathfrak{gl}_3$ versus $\mathfrak{sl}_3$}: projecting the weights to the
$\mathfrak{sl}_3$ weight lattice turns $s_\lambda$ into the irreducible character
$\chi_{\bar\lambda}$, and $\ch B_7=\chi_{(1,1)}-\chi_{(0,0)}$ (in fundamental-weight
coordinates); the coefficient of the trivial character is $-1$.
\item \emph{Polynomial versus Laurent}: all weights of $B_7$ are nonnegative, so $\ch B_7$ is
an honest polynomial and the expansion is in ordinary Schur polynomials.
\end{itemize}

\section{Machine verification}

\paragraph{Lean 4} (\texttt{lean4/Main.lean}, Lean 4.19.0, core library only, no \texttt{sorry},
no \texttt{native\_decide}; the main theorems use only \texttt{propext} and \texttt{Quot.sound},
some lemmas also \texttt{Classical.choice}). Everything is defined from scratch:
\begin{itemize}
\item \texttt{Crystal n}: $\wt\colon\texttt{Fin n}\to\Z^3$, $\tilde e_i,\tilde f_i\colon
\texttt{Fin n}\to\texttt{Option (Fin n)}$, $\varepsilon_i,\varphi_i\colon\texttt{Fin n}\to\Z$;
\texttt{IsCrystal} (axioms 1--4) and \texttt{IsSeminormal}; \texttt{seminormal\_eps\_max} and
\texttt{seminormal\_phi\_max} prove that it means $\varepsilon_i=\max\{k:\tilde e_i^kb\neq0\}$
(and likewise for $\varphi_i$).
\item Laurent polynomials as finite formal sums with coefficient-wise equality \texttt{PolyEq}
(a sound and complete decision procedure is proved), multiplication, the $S_3$-action, $J$, the
orbit sum, the character \texttt{ch}.
\item \texttt{schur}: Schur polynomials as sums over semistandard tableaux (all fillings of the
shape with entries in $\{1,2,3\}$ are generated and filtered); the bialternant identities
$s_\lambda a_\rho=a_{\lambda+\rho}$ for $\lambda\vdash3$; \texttt{schur3\_indep} (uniqueness of
degree-$3$ Schur expansions).
\item \texttt{WeylSym f q}: $q\cdot J(x^\rho)=J(x^\rho f)$; \texttt{schurCoeffW}:
$[x^{\lambda+\rho}]J(x^\rho f)$.
\item \texttt{B7\_isCrystal}, \texttt{B7\_seminormal}, \texttt{ch\_B7\_expands}
($\ch B_7=s_{210}-s_{111}$), \texttt{weylSym\_B7}, \texttt{weyl\_forced} (for \emph{every} $q$
with \texttt{WeylSym (ch B7) q} and every degree-3 expansion of $q$ the coefficients are
$(0,1,-1)$), \texttt{orbit\_forced}, \texttt{avg\_forced}, \texttt{B7\_counterexample}.
\item Main theorems: \texttt{conjecture\_00000003831\_false} : $\neg$\,\texttt{ClaimWeyl} and the variants
\texttt{\_abstract}, \texttt{\_bialt}, \texttt{\_orbit}, \texttt{\_avg} for the other readings.
\item Non-vacuity: \texttt{B100\_ok} ($B(1,0,0)$ is a seminormal crystal with character
$s_{100}$), \texttt{B210\_ok} ($B(2,1,0)$ satisfies all hypotheses and has expansion $s_{210}$,
nonnegative), \texttt{broken\_rejected} (a datum violating axiom~1 is rejected).
\end{itemize}
Limitations: in Lean the Schur expansions quantified over are the degree-$3$ ones
$a s_{300}+b s_{210}+c s_{111}$ (sufficient because $\ch B_7$ is homogeneous of degree~3; general
uniqueness is Lemma~\ref{lem:unique}). That the symmetrizer fixes symmetric polynomials is checked
for $\ch B_7$, not proved in general. Normality, Lemma~\ref{lem:semi}, the $\mathfrak{gl}_2$
statement and minimality are in this report and \texttt{verify.py} only.

\paragraph{Python} (\texttt{verify.py}, standard library). It checks the axioms by a different
method (string lengths computed by walking), computes Schur polynomials both from tableaux and
by exact division $a_{\lambda+\rho}/a_\rho$, computes the symmetrizer by exact polynomial
division and via Demazure operators $\pi_1\pi_2\pi_1$, the orbit sum and average (with
fractions), runs the $\mathfrak{gl}_2$ (up to $10$ vertices) and $\mathfrak{gl}_3$ (up to $7$
vertices) searches of Proposition~\ref{prop:min}, and sanity-checks the normal crystals
$B(1,0,0)$, $B(2,1,0)$ and $B(1)^{\otimes3}$ (tensor rule).

\begin{thebibliography}{9}
\bibitem{BS17} D.~Bump and A.~Schilling, \emph{Crystal Bases: Representations and
Combinatorics}, World Scientific, 2017.
\bibitem{Kas93} M.~Kashiwara, The crystal base and Littelmann's refined Demazure character
formula, \emph{Duke Math.\ J.} 71 (1993), 839--858.
\bibitem{Kas95} M.~Kashiwara, On crystal bases, in: \emph{Representations of Groups}
(Banff, 1994), CMS Conf.\ Proc.\ 16, Amer.\ Math.\ Soc., 1995, 155--197.
\bibitem{Ste03} J.~R.~Stembridge, A local characterization of simply-laced crystals,
\emph{Trans.\ Amer.\ Math.\ Soc.} 355 (2003), 4807--4823.
\end{thebibliography}

\end{document}
